% Options for packages loaded elsewhere
\PassOptionsToPackage{unicode}{hyperref}
\PassOptionsToPackage{hyphens}{url}
%
\documentclass[
  11pt,
  a4paper,
]{article}
\usepackage{amsmath,amssymb}
\usepackage{iftex}
\ifPDFTeX
  \usepackage[T1]{fontenc}
  \usepackage[utf8]{inputenc}
  \usepackage{textcomp} % provide euro and other symbols
\else % if luatex or xetex
  \usepackage{unicode-math} % this also loads fontspec
  \defaultfontfeatures{Scale=MatchLowercase}
  \defaultfontfeatures[\rmfamily]{Ligatures=TeX,Scale=1}
\fi
\usepackage{lmodern}
\ifPDFTeX\else
  % xetex/luatex font selection
  \setmainfont[]{DejaVu Serif}
  \setsansfont[]{DejaVu Sans}
\fi
% Use upquote if available, for straight quotes in verbatim environments
\IfFileExists{upquote.sty}{\usepackage{upquote}}{}
\IfFileExists{microtype.sty}{% use microtype if available
  \usepackage[]{microtype}
  \UseMicrotypeSet[protrusion]{basicmath} % disable protrusion for tt fonts
}{}
\makeatletter
\@ifundefined{KOMAClassName}{% if non-KOMA class
  \IfFileExists{parskip.sty}{%
    \usepackage{parskip}
  }{% else
    \setlength{\parindent}{0pt}
    \setlength{\parskip}{6pt plus 2pt minus 1pt}}
}{% if KOMA class
  \KOMAoptions{parskip=half}}
\makeatother
\usepackage{xcolor}
\usepackage[margin=20mm]{geometry}
\setlength{\emergencystretch}{3em} % prevent overfull lines
\providecommand{\tightlist}{%
  \setlength{\itemsep}{0pt}\setlength{\parskip}{0pt}}
\setcounter{secnumdepth}{-\maxdimen} % remove section numbering
\usepackage{setspace}
\usepackage{lineno}
\usepackage{fancyhdr}
\usepackage{sectsty}
\allsectionsfont{\raggedright}
\usepackage{float}
\floatplacement{figure}{H}
\setlength{\emergencystretch}{3em}
\widowpenalty=10000
\clubpenalty=10000
\setlength{\tabcolsep}{3pt}
\pagestyle{fancy}\fancyhf{}
\fancyhead[L]{\footnotesize Supplementary Information | Tidal potential measured with distant optical clocks}
\fancyfoot[C]{\thepage}
\setlength{\headheight}{14pt}
\ifLuaTeX
  \usepackage{selnolig}  % disable illegal ligatures
\fi
\IfFileExists{bookmark.sty}{\usepackage{bookmark}}{\usepackage{hyperref}}
\IfFileExists{xurl.sty}{\usepackage{xurl}}{} % add URL line breaks if available
\urlstyle{same}
\hypersetup{
  hidelinks,
  pdfcreator={LaTeX via pandoc}}

\author{}
\date{}

\begin{document}

\onehalfspacing
\linenumbers

\hypertarget{supplementary-information}{%
\section{Supplementary Information}\label{supplementary-information}}

\textbf{Tidal potential measured with distant optical clocks}

This document accompanies the main article. References, figure numbers,
table numbers and named Methods sections follow the main article.

\hypertarget{supplementary-methods}{%
\subsection{Supplementary Methods}\label{supplementary-methods}}

\hypertarget{sm1.-physical-model-implementation-and-sensitivity}{%
\subsubsection{SM1. Physical model implementation and
sensitivity}\label{sm1.-physical-model-implementation-and-sensitivity}}

Here \(V_{2m}\) are the degree-2 tide-generating potentials of the Moon
and Sun separated by order \(m\), \(V_3\) is the lunar degree-3
potential, and \(u_r\) is radial displacement. We use
\(k_{20}=0.29525\), \(k_{22}=0.29801\) and \(k_3=0.093\). The IERS 2010
resonance formula is applied constituent by constituent for \(k_{21}\).
It gives 0.2588 for the lunisolar diurnal tide, K₁, 0.2873 for the
principal solar diurnal tide, P₁, and 0.2975 for the principal lunar
diurnal tide, O₁. These diurnal constituents describe approximately
once-daily changes associated with lunar and solar declination. The
radial displacement is computed with the IERS algorithm, including the
frequency and latitude dependence of the Love number \(h\)
\textsuperscript{20}.

The ocean elevation \(\zeta\) is taken from FES2022 on a \(1/30^\circ\)
grid \textsuperscript{21,22}, \(\rho_w=1,025\ \mathrm{kg\,m^{-3}}\),
\(a\) and \(M_E\) are Earth's radius and mass, and \(\psi\) is angular
distance. The Green's function follows ref. \textsuperscript{23}. PREM
load Love numbers are used through degree \(n=696\)
\textsuperscript{29}, from the gravity-toolkit distribution
\textsuperscript{30}. Above this degree, \(h'_n\) is joined smoothly to
its Boussinesq limit of −6.22, and the series is summed with a Kummer
transformation through \(n=40,000\). The convolution uses full grid
resolution within 3° of each station and sums blocks of \(8\times8\)
cells farther away. This aggregation changes differential M₂ and K₁
loading by at most 0.05\%. The ocean mean is removed to enforce mass
conservation. The larger lunar elliptic diurnal tide, Q₁, accounts for a
further once-daily contribution associated with the Moon's elliptical
orbit. These four diurnal constituents and the four semidiurnal
constituents introduced above (N₂, M₂, S₂ and K₂) are synthesized with
astronomical arguments and nodal corrections.

The full loading potential exceeds the vertical-displacement
contribution \(-g\delta h\) by factors of approximately 1.4 for the
semidiurnal constituents and 1.8 for the diurnal constituents. The
radial displacement from the same convolution agrees with the FES2022
load-tide grid to within 2\% in amplitude and 1.2° in phase for seven of
the eight constituents. This comparison checks the convolution and
displacement contribution; the direct attraction and induced-potential
terms use the same PREM load Love numbers. For S₂, the amplitude differs
by 18\% and the phase at Wuhan by 54°; the load-tide grid uses a
different ocean-model release. The S₂ loading contributions expressed as
fractional frequency are \(0.62\times10^{-18}\) in Shanghai and
\(0.18\times10^{-18}\) in Wuhan, giving a differential amplitude of
\(0.63\times10^{-18}\), compared with \(1.42\times10^{-18}\) for M₂.
Substituting the S₂ loading inferred from the load-tide grid changes the
full-template amplitude by 0.001 and the semidiurnal amplitude by 0.000
at the reported precision, with no reported change in the lunar
amplitude. The lunar and near-12-h groups have separate fitted
amplitudes.

The Shanghai clock is 22 km from the nearest ocean cell. To examine
coastal-mask sensitivity, land cells within 0.3° or 0.6° of the
coastline in a \(4^\circ\times5^\circ\) box centred on Shanghai are
filled with the nearest ocean values. These strips, approximately 33--66
km wide, increase the Shanghai M₂ loading potential expressed as
\(\delta W/g\) from 15.3 mm to 17.0 or 18.5 mm and decrease the fitted
response by 0.06 or 0.07. These are physical-model sensitivity tests,
with shifts approximately half the quoted lunar statistical uncertainty;
they are not treated as independent Gaussian errors or as an exhaustive
loading-model uncertainty.

The eight constituents explain 99.84\% of the model variance. The
remaining component, with an r.m.s. of \(0.18\times10^{-18}\), is
retained at its predicted value. Setting it to zero or doubling it gives
semidiurnal amplitudes of 0.95--0.96 and lunar amplitudes of 0.95--0.99.

\hypertarget{sm2.-selection-and-timing-checks}{%
\subsubsection{SM2. Selection and timing
checks}\label{sm2.-selection-and-timing-checks}}

The transient criterion uses no tidal model. Adding the model with
factors of −1, 1 or 3 leaves the removed samples unchanged. Six robust
standard deviations is the lowest tested threshold for which this
invariance holds up to three times the model. For block lengths of 60 or
300 s and thresholds from four to eight standard deviations, the lunar
amplitude ranges from 0.93 to 1.02. Retaining the excursions gives
\(0.96\pm0.14\); the full-template and semidiurnal amplitudes change by
less than 0.04. Before transient removal, the procedure reproduces the
summer sample set of ref. \textsuperscript{25}. After removal, the
campaign diurnal amplitudes are mutually consistent (\(p=0.21\)).
Because the tested tidal additions do not change selection and the
subsequent filter is linear, the injection simulations use the selected
data.

Durations refer to retained 1-s samples; exposure measured from the
filtered data at 600-s steps, used for temperature overlap, is slightly
shorter. The last April segment ends at 18:00 on 19 April, following the
experiment team's segment table. Omitting endpoint trimming while
retaining the other steps changes the lunar amplitude by −0.009 and the
full-template amplitude by 0.008. Each segment has its own fitted
offset. Time stamps are in local time (UTC+8); shifting the model by ±10
min changes \(A\) by less than 0.02. Long-period tides are outside the
analysis: the similar site latitudes make the predicted differential
semiannual tide only about \(2\times10^{-20}\).

\hypertarget{sm3.-primary-uncertainty-calibration-and-injection-validation}{%
\subsubsection{SM3. Primary uncertainty calibration and injection
validation}\label{sm3.-primary-uncertainty-calibration-and-injection-validation}}

For the lunar group, two independent sets of 200 realizations give
standard deviations of 0.13 and 0.14, compared with 0.13 for the larger
of the scaled least-squares and jackknife errors. We quote 0.14; the
resampling bias is +0.01. With the daily disturbances defined in the
main Methods, 5,000 realizations per scenario give standard deviations
rounding to 0.14, coverage of 0.67--0.68 for the interval ±0.14, and
0.945--0.952 for ±1.96 × 0.14.

For the semidiurnal response, the group-bootstrap standard deviation of
0.08 is smaller than the calibrated value of 0.12, and the larger value
is quoted. Segment-bootstrap standard deviations give a lunar amplitude
of \(0.98\pm0.11\) and phase of \(-6^\circ\pm8^\circ\).

The added disturbances have r.m.s. values up to \(1.9\times10^{-18}\).
In 300 realizations per scenario and injected amplitude, lunar-amplitude
biases lie between −0.03 and −0.00 at the reported precision. The larger
of the scaled least-squares and jackknife errors covers the injected
value in 0.59--0.71 of realizations at one standard error and 0.90--0.95
at 1.96 standard errors. The main estimate uses the calibrated
uncertainty above. For \(A=1\), the fitted lunar amplitude exceeds three
times its estimated error in every realization; for \(A=0\), this
fraction is at most 0.003. These simulation results are conditional on
the selected residual and disturbance models.

In the triangular-filter reduction, the correlation between the lunar
and neighbouring semidiurnal response coefficients is −0.32. The
full-template and semidiurnal responses are \(0.90\pm0.12\) and
\(0.96\pm0.12\), respectively (Extended Data Fig. 11). Summer supplies
61\% of the conditional information on the lunar response and gives
\(0.87\pm0.18\), compared with \(1.05\pm0.31\) for the earlier campaigns
together. The five lunar estimates are consistent (\(\chi^2=4.1\), four
degrees of freedom, \(p=0.40\); Extended Data Table 3). Semidiurnal
estimates also agree across campaigns (\(\chi^2=3.2\), four degrees of
freedom, \(p=0.52\); Extended Data Fig. 3). Resampling the operating
groups gives a 95\% lunar interval of {[}0.76, 1.18{]}.

\hypertarget{sm4.-primary-prediction-monitoring-and-environmental-checks}{%
\subsubsection{SM4. Primary prediction, monitoring and environmental
checks}\label{sm4.-primary-prediction-monitoring-and-environmental-checks}}

Held-out prediction estimates the amplitude from the other observation
periods and applies it to the excluded period, fitting that period's
segment offsets. The whole-template prediction reduces residual sums of
squares in seven of eight periods, by 4.9\% when pooled; the lunar
prediction improves six of eight periods, by 1.7\% when pooled. These
checks use the same underlying record as the primary fits.

Over the 98-h overlap, the blackbody shift calculated from chamber
temperature and filtered like the ratio has an r.m.s. of
\(2.8\times10^{-18}\), including the 30 June transient. Adding it with a
free coefficient gives \(0.92\pm0.74\), or \(0.72\pm0.74\) after
excluding the transient, with day-by-day jackknife errors. Both are
consistent with zero at this precision. Through the summer
lunar-estimator response of 0.016 per \(10^{-18}\), a residual with the
size of the observed 24-h ratio term would change the lunar amplitude by
about 0.045. Without real-time correction, the blackbody shift would
contain a daily component of \(2.0\times10^{-18}\) peaking at 12:13
local time. In the same overlap, the fitted 24-h ratio term is
\((0.9\pm1.1)\times10^{-18}\) with maximum at 10:54 ± 5.0 h. The mean
daily profile of the uncorrected blackbody contribution has a
peak-to-peak range of \(4\times10^{-18}\), indicating the size of the
removed effect.

On the primary sampling after segment-mean removal, the predicted
solid-tide and loading root-mean-square fractional frequencies are
\(3.0\times10^{-18}\) and \(1.35\times10^{-18}\), and their sum is
\(4.1\times10^{-18}\). The whole-template response fitted to the
solid-Earth tide alone is \(1.19\pm0.20\), compared with 1.32 expected
on this sampling for the combined model. The separately fitted loading
coefficient is \(0.89\pm0.55\), consistent with the combined prediction
but insufficient to establish a separate loading detection.

The summer 24-h term has amplitude \((2.8\pm1.3)\times10^{-18}\) and
peaks at 15:48 local time, with a 1.5-h standard uncertainty; the 12-h
term is \((0.6\pm0.7)\times10^{-18}\). Their summer projection
coefficients are 0.016 and 0.043 per \(10^{-18}\), giving maximum shifts
of 0.045 and 0.024 at those fitted amplitudes (Extended Data Table 3).
For the full record, a common-phase 12-h perturbation of amplitude
\(10^{-18}\) changes the whole-template response by at most 0.087.

Campaign-specific 24-h regressors with phases fixed by local air
temperature give a lunar response of \(0.99\pm0.14\); free-phase
campaign-specific 24-h terms give \(0.97\pm0.15\). Omitting each earlier
campaign in turn gives lunar responses between 0.89 and 0.99. Three of
the nine identified transient excursions coincide with approximately
equal and opposite link-monitor excursions (Extended Data Fig. 5). The
transient criterion acts on 300-s deviations from a local running median
and does not use the tidal waveform.

Temperature-anomaly regressors at both sites give a lunar response of
\(0.98\pm0.13\). Replacing segment offsets by operating-group or period
offsets gives \(0.97\pm0.11\) and \(0.99\pm0.11\), respectively. These
estimates use the triangular-filter reduction (Fig. 4a).

In the triangular-filter reduction, fits with common 24-h and 12-h terms
give a lunar response of \(1.00\pm0.13\). Campaign-specific free-phase
12-h terms give \(0.82\pm0.17\), or \(0.86\pm0.18\) when 24-h terms are
also included (Fig. 4). Their errors follow the
larger-of-scaled-least-squares-and-jackknife prescription for fits with
empirical daily terms.

\hypertarget{sm5.-primary-spectral-displays-and-alternative-processing}{%
\subsubsection{SM5. Primary spectral displays and alternative
processing}\label{sm5.-primary-spectral-displays-and-alternative-processing}}

As a processing cross-check, the segment-wise estimator of ref.
\textsuperscript{25} demeans filtered data and model within each
segment, fits the amplitude through the origin using non-overlapping
1,200-s windows, and combines segment estimates with inverse-variance
weights. It differs from the primary joint fit in both sampling and
weighting. For the summer record it gives \(A=0.78\pm0.19\) with the
present physical model, compared with \(0.87\pm0.15\) with the model of
ref. \textsuperscript{25}; the primary joint lunar fit gives
\(0.87\pm0.18\). The lower summer whole-template response,
\(0.67\pm0.18\), accompanies a low diurnal response and should not be
substituted for the lunar estimate. For April, the segment-wise and
joint estimates are \(0.66\pm0.31\) and 0.91. For the complete record,
the segment-wise result is \(0.84\pm0.12\), with uncertainty scaled by
the square root of the segment-to-segment reduced chi-square of 2.1.

Independent processing of the same 2026 records uses non-overlapping
600-s means, corrected clock systematic shifts and fixed ten-day
calendar baselines. It gives a whole-template response of
\(0.88\pm0.10\) without empirical daily terms and a lunar response of
\(0.89\pm0.28\) with daily terms in each ten-day block. Its
whole-template prediction improves all six held-out 2026 periods.
Simulations with white, Ornstein--Uhlenbeck (correlation times of 1 h to
3 d), Student-\(t\) and moving-average noise give biases below 0.04 and
nominal 95\% interval coverage of 94.5--99.8\%, using the largest of
least-squares, delete-day and delete-period standard errors. Simpler
solid-tide approximations change its whole-template scale by less than
0.01. These estimates are processing checks on shared observations and
are not combined as independent measurements.

Extended Data Fig. 11 collects the triangular-filter grouped responses,
their covariance and the full-record harmonic spectrum. In the
three-group fit, the diurnal coefficient is \(0.35\pm0.51\); the
two-band estimate in Extended Data Fig. 3b is \(0.36\pm0.50\). The
sampling, physical templates and uncertainty procedures are defined in
``Tidal fits and uncertainty''.

The full-record spectrum displays semidiurnal structure that is strongly
reduced after subtraction of the fitted whole tidal template (Extended
Data Fig. 11c). At M₂, harmonic power after projection of segment
offsets decreases by 99\%. The April time series provides a
complementary view on an actual calendar-time axis (Extended Data Fig.
4). These spectra show squared sinusoidal amplitude divided by two on a
linear ordinate, rather than density per unit frequency. Their
residual-based thresholds provide frequency-domain evidence alongside
the physical-template estimates; the power reduction is not an
instrumental signal-to-noise ratio.

\hypertarget{sm6.-auxiliary-response-and-uncertainty-details}{%
\subsubsection{SM6. Auxiliary response and uncertainty
details}\label{sm6.-auxiliary-response-and-uncertainty-details}}

The full-template estimate is \(0.9551\pm0.1082\), and the grouped lunar
and neighbouring estimates are \(1.0004\pm0.1320\) and
\(0.9251\pm0.3649\). In the phase-relaxed fit, the lunar in-phase and
quadrature estimates are \(1.0007\pm0.1320\) and \(-0.0969\pm0.1989\).

The full and lunar responses are close to the common-white-noise
expectation; the other empirical errors are larger. Loading is weakly
constrained, with projected solid-Earth/loading template correlation
0.725 and coefficient correlation −0.888 under day deletion. This
comparison does not determine a loading coherence exponent. For the
triangular-filter results, the corresponding ratios are 6.93 for the
lunar group, 3.10 for the neighbouring semidiurnal group and 0.72 for
the two-band diurnal estimate. These ratios retain that reduction's own
uncertainty definitions.

For M₂ and S₂, at the measured amplitudes \(4.5047\times10^{-18}\) and
\(1.9214\times10^{-18}\), the observing-period noise propagation gives
local standard errors \(0.6153\times10^{-18}\) and
\(0.6788\times10^{-18}\).

For the diurnal-group and auxiliary loading coefficients, respectively,
the period-noise standard errors are 0.3631 and 0.4337; their
UTC-day-jackknife errors are 0.5381 and 0.5438.

As an additional noise comparison, we propagate the measured
observing-period noise coefficients through the unchanged estimators
using \(C_\beta=H\Sigma H^{\mathsf T}\) and
\(\Sigma_{kk}=\alpha_{p(k)}^2/\Delta t_k\). Like the common-\(\alpha\)
calculation, this assumes independent white-frequency-noise bin errors;
``period'' identifies the observing period and does not denote a
periodic disturbance. The resulting \(S/R\) values for M₂, S₂, the
diurnal group and ocean loading are 7.32, 2.83, 1.59 and 1.12. These
supplementary comparisons do not change the weights or measured
amplitudes of the estimators used in Table 1.

\hypertarget{sm7.-compensation-alias-and-correlation-diagnostics}{%
\subsubsection{SM7. Compensation, alias and correlation
diagnostics}\label{sm7.-compensation-alias-and-correlation-diagnostics}}

For donor \(D\) and target \(T\), the joint fit supplied the fixed donor
coefficient vector \(\hat\beta_D\). For 61 values of λ from 0 to 1.5, we
formed \(y_\lambda=y-\lambda X_D\hat\beta_D\) and fitted only the target
pair together with \(N\). With \(P_T=R_NX_T\) and \(P_D=R_NX_D\), its
coefficient vector obeys

\[
\hat\beta_T(\lambda)=\hat\beta_{T,\mathrm{single}}-\lambda P_T^+P_D\hat\beta_D.
\]

Here \(P_T^+\) denotes the Moore--Penrose pseudoinverse. We also
repeated the sweep using the donor's single-frequency estimate to show
the dependence on waveform definition. A separate jointly controlled fit
retains both M₂ and S₂ pairs after subtraction. Its target coefficient
is necessarily invariant, because the subtracted waveform lies within
the fitted donor subspace; this curve is an algebraic control, not
evidence against leakage. Amplitude is the norm of the two coefficients
and phase is \(\operatorname{atan2}(\beta_{\sin},\beta_{\cos})\),
corresponding to \(A\cos(2\pi f(t-t_0)-\phi)\). Target-only incremental
improvement is
\(\Delta\mathrm{RSS}=\hat\beta_T^\mathsf TP_T^\mathsf TP_T\hat\beta_T\).
The jointly controlled improvement instead residualizes the target
against both \(N\) and the donor, so these two improvements have
different baselines even where their coefficient estimates coincide.
Density panels use the normalization defined in ``Auxiliary
frequency-domain analysis''.

Canonical correlations are the singular values of \(Q_M^\mathsf TQ_S\),
where each \(Q\) spans a residualized constituent pair. Template lag
correlations evaluate the fitted donor waveform at \(t-\ell\) and the
target at \(t\), residualize both on the unchanged observation support,
and normalize their inner product. Lags span ±24 h in 5-min steps. Their
zero-lag correlation is 0.1955; maximum absolute lag correlations are
0.2413 for shifted M₂ against S₂ and 0.2113 in the reverse direction.
These are design and fitted-template similarities, not an observed
stochastic cross-correlation function or a measured propagation delay.

For alias diagnostics, we scanned 20--25 μHz with spacing
\(1/(12T_{\rm span})\), where \(T_{\rm span}\) is the elapsed
observation span. For each donor, we selected the three strongest local
maxima of its pure-injection spectrum at least \(2/T_{\rm span}\) away
from its own frequency. Selection used the injection response, not a
significance threshold on the observed data. At each selected frequency,
we compared the actual refit to the predicted vector response
\(\hat b_f(\lambda)=\hat b_f(0)-\lambda\hat b_{f,D}\). Its squared norm
in the frequency-specific Gram metric gives the predicted quadratic
power response. Numerical agreement verifies the processing operator;
because this is a linear identity, agreement alone cannot establish an
exclusively aliased origin of the observed peak. The compensation curves
describe deterministic sensitivity conditional on the fitted waveforms.
In this non-overlapping reduction, the lunar-waveform compensation
minima at the sidelobes near 22.1016 and 22.6283 μHz occur at 0.996 and
1.0445 times the jointly fitted waveform (Extended Data Fig. 10).

We express each four-coefficient covariance as a correlation matrix,
\(R_{ij}=C_{ij}/\sqrt{C_{ii}C_{jj}}\), separately for scaled least
squares and UTC-day deletion (Extended Data Fig. 9c,d). These matrices
describe the coupled uncertainty of the joint coefficient estimates,
distinct from correlations of the design columns and fitted templates.

Identification of the S₂-frequency contribution does not by itself
assign its entire amplitude to the astronomical tide: a laboratory
variation at exactly 12 h occupies the same frequency. Environmental and
daily-term checks are described above.

The largest absolute cross-component coefficient correlations between
the joint M₂ and S₂ estimates are 0.248 under scaled least squares and
0.261 under UTC-day deletion (Extended Data Fig. 9). These coefficient
correlations describe uncertainty coupling and differ from the canonical
correlation of the projected design pairs.

At 22.1016 and 22.6283 μHz, the lunar sampling response matches features
whose powers fall by 99.56\% and 99.39\% after lunar subtraction
(Extended Data Fig. 10).

\hypertarget{references}{%
\subsection{References}\label{references}}

Reference numbers are retained from the main article.

\textbf{20.} Petit, G. \& Luzum, B. (eds) \emph{IERS Conventions
(2010)}. IERS Technical Note \textbf{36} (2010).
\url{https://iers-conventions.obspm.fr/content/tn36.pdf}

\textbf{21.} Lyard, F. H., Allain, D. J., Cancet, M., Carrère, L. \&
Picot, N. FES2014 global ocean tide atlas: design and performance.
\emph{Ocean Sci.} \textbf{17}, 615--649 (2021).
\url{https://doi.org/10.5194/os-17-615-2021}

\textbf{22.} CNES. \emph{FES2022 (Finite Element Solution) Tidal model
(Version 2024)} {[}Data set{]} (2024).
\url{https://doi.org/10.24400/527896/A01-2024.004}

\textbf{23.} Farrell, W. E. Deformation of the Earth by surface loads.
\emph{Rev.~Geophys.} \textbf{10}, 761--797 (1972).
\url{https://doi.org/10.1029/RG010i003p00761}

\textbf{25.} Hou, L. et al.~Yb/Sr optical clock comparison at
\(2.2\times10^{-18}\) uncertainty over a 1350-km fiber link. Preprint at
\url{https://arxiv.org/abs/2610.34250v1} (2026).

\textbf{29.} Han, D. \& Wahr, J. The viscoelastic relaxation of a
realistically stratified earth, and a further analysis of postglacial
rebound. \emph{Geophys. J. Int.} \textbf{120}, 287--311 (1995).
\url{https://doi.org/10.1111/j.1365-246x.1995.tb01819.x}

\textbf{30.} Sutterley, T. C. \emph{gravity-toolkit: Python tools for
working with time-variable gravity fields}.
\url{https://github.com/polargeodesy/gravity-toolkit}.

\end{document}
